% figures/mechanism.pdf -- how the anode degrades the molecule.
%
% Water is oxidised at the anode surface to a sorbed hydroxyl radical, which then attacks the
% aromatic ring of the pollutant. Only the drawing lives here: LabHarness adds the journal
% style around it, which is why there is no preamble and no font in this file.
%
% electron pair   -- a full arrowhead, for a pair of electrons
% single electron -- a half arrowhead, for a single electron
\begin{tikzpicture}[node distance=1.1cm]
  \node (water) {\chemfig{H-O-H}};
  \node (anode) [right=of water] {M};
  \node (radical) [right=2.6cm of anode] {\chemfig{M\cdot OH}};
  \node (ring) [right=2.4cm of radical] {\chemfig{*6(-=-=(-R)-=)}};

  \draw[electron pair] ($(anode.east)+(0.2,0)$) -- node[node label, above] {$-$H$^+$, $-$e$^-$}
    ($(radical.west)+(-0.2,0)$);
  \draw[single electron] ($(radical.east)+(0.2,0)$) -- node[node label, above] {$\cdot$OH}
    node[node label, below] {attack} ($(ring.west)+(-0.2,0)$);
\end{tikzpicture}
